% =====================================================================
%  preamble.tex —— Orlov《凝聚层的导出范畴及其之间的等价》中文译注本
%  导言区 · Lead 独占，译者只读
%  宏包与骨架沿用本库房规（SY08-01 体例），仅替换与本稿相关的部分
% =====================================================================
\usepackage[b5paper,top=21mm,bottom=21mm,left=19mm,right=19mm,headheight=14pt,headsep=5mm,footskip=9mm]{geometry}
\usepackage{amsmath,amssymb,mathrsfs,bm}
\usepackage{graphicx}
\usepackage{xcolor}
\usepackage{booktabs,array,multirow,longtable,multicol}
\usepackage{newunicodechar}% 兜底：正文偶发的 Unicode 符号也能排版
\newunicodechar{■}{\ensuremath{\blacksquare}}
\newunicodechar{□}{\ensuremath{\square}}
\newunicodechar{≥}{\ensuremath{\geq}}
\newunicodechar{≤}{\ensuremath{\leq}}
\newunicodechar{≠}{\ensuremath{\neq}}
\newunicodechar{≡}{\ensuremath{\equiv}}
\newunicodechar{≅}{\ensuremath{\cong}}
\newunicodechar{∼}{\ensuremath{\sim}}
\newunicodechar{≈}{\ensuremath{\approx}}
\newunicodechar{≃}{\ensuremath{\simeq}}
\newunicodechar{±}{\ensuremath{\pm}}
\newunicodechar{×}{\ensuremath{\times}}
\newunicodechar{⋅}{\ensuremath{\cdot}}
\newunicodechar{⊗}{\ensuremath{\otimes}}
\newunicodechar{⊕}{\ensuremath{\oplus}}
\newunicodechar{⊠}{\ensuremath{\boxtimes}}
\newunicodechar{⌈}{\ensuremath{\lceil}}
\newunicodechar{⌉}{\ensuremath{\rceil}}
\newunicodechar{⌊}{\ensuremath{\lfloor}}
\newunicodechar{⌋}{\ensuremath{\rfloor}}
\newunicodechar{∞}{\ensuremath{\infty}}
\newunicodechar{→}{\ensuremath{\to}}
\newunicodechar{↦}{\ensuremath{\mapsto}}
\newunicodechar{↪}{\ensuremath{\hookrightarrow}}
\newunicodechar{⊂}{\ensuremath{\subset}}
\newunicodechar{⊆}{\ensuremath{\subseteq}}
\newunicodechar{∈}{\ensuremath{\in}}
\newunicodechar{≺}{\ensuremath{\prec}}
\newunicodechar{⩽}{\ensuremath{\leqslant}}
\newunicodechar{Δ}{\ensuremath{\Delta}}
\newunicodechar{Σ}{\ensuremath{\Sigma}}
\newunicodechar{Γ}{\ensuremath{\Gamma}}
\newunicodechar{Λ}{\ensuremath{\Lambda}}
\newunicodechar{Φ}{\ensuremath{\Phi}}
\newunicodechar{Ψ}{\ensuremath{\Psi}}
\newunicodechar{ξ}{\ensuremath{\xi}}
\newunicodechar{η}{\ensuremath{\eta}}
\newunicodechar{λ}{\ensuremath{\lambda}}
\newunicodechar{μ}{\ensuremath{\mu}}
\newunicodechar{ν}{\ensuremath{\nu}}
\newunicodechar{ω}{\ensuremath{\omega}}
\newunicodechar{φ}{\ensuremath{\varphi}}
\newunicodechar{ψ}{\ensuremath{\psi}}
\newunicodechar{・}{\textperiodcentered}
% 西里尔字母（俄文书名）：latin modern 无西里尔字形，借系统 Times New Roman 兜底
\IfFontExistsTF{Times New Roman}
  {\newfontfamily\cyrfont{Times New Roman}}
  {\IfFontExistsTF{DejaVu Serif}
     {\newfontfamily\cyrfont{DejaVu Serif}}
     {\newcommand{\cyrfont}{}}}
\usepackage{enumitem}
\usepackage{fancyhdr}
\usepackage{caption}
\usepackage{tcolorbox}
\tcbuselibrary{most}
\usepackage{tikz}
\usepackage{tikz-cd}
\usepackage{hyperref}
\hypersetup{colorlinks=true,linkcolor=black,citecolor=black,urlcolor=blue!45!black,
            bookmarksnumbered=true,
            pdftitle={凝聚层的导出范畴及其之间的等价（Orlov 著，中文译注本）},
            pdfauthor={D. O. Orlov}}

\setlength{\parindent}{2em}
\linespread{1.12}
\allowdisplaybreaks[2]
\setlength{\emergencystretch}{6em}
\tolerance=3000
\hbadness=10000

% ---------- 页眉页脚 ----------
\pagestyle{fancy}
\fancyhf{}
\renewcommand{\headrulewidth}{0.3pt}
\fancyhead[LE]{\small\thepage}
\fancyhead[RE]{\small\leftmark}
\fancyhead[RO]{\small\thepage}
\fancyhead[LO]{\small\rightmark}
\fancyfoot{}
\renewcommand{\chaptermark}[1]{\markboth{\CTEXthechapter\quad #1}{}}
\renewcommand{\sectionmark}[1]{\markright{\thesection\quad #1}}

% ---------- 章节标题格式 ----------
\ctexset{
  chapter = {
    format = \centering\Large\bfseries,
    name = {第,章},
    number = \arabic{chapter},
    beforeskip = 18pt,
    afterskip = 22pt,
  },
  section = {
    format = \large\bfseries,
    beforeskip = 14pt,
    afterskip = 10pt,
  },
  subsection = {
    format = \normalsize\bfseries,
    beforeskip = 11pt,
    afterskip = 6pt,
  },
}

% ---------- 公式编号 ----------
% 原书公式按章连续编号 (1)–(31)，跨章不重起。此处一律手工 \tag{原编号}，
% 故不安排 equation 自动计数；\numberwithin 仅作兜底。
\numberwithin{equation}{chapter}

% ---------- 陈述环境 ----------
% 原书（RMS 2003 英文版）定理类陈述末尾**不带任何符号**，证明末尾亦不带；
% 本译本照此不加尾符，以存原貌。
% 用法：\begin{theorem}{2.1.3} 正文…… \end{theorem}
% 可选参数 [附加名]：\begin{theorem}[Orlov]{3.2.1}
% 编号是**必选第一参数**，一律照抄原书编号（形如 1.1.1、2.2.8），不用自动编号。
\newcommand{\fkstmt}[3][]{\par\medskip\noindent\textbf{#2~#3}%
  \if\relax\detokenize{#1}\relax\else\textbf{（#1）}\fi\ \ignorespaces}
\newenvironment{theorem}[2][]{\fkstmt[#1]{定理}{#2}}{\par\medskip}
\newenvironment{proposition}[2][]{\fkstmt[#1]{命题}{#2}}{\par\medskip}
\newenvironment{lemma}[2][]{\fkstmt[#1]{引理}{#2}}{\par\medskip}
\newenvironment{corollary}[2][]{\fkstmt[#1]{推论}{#2}}{\par\medskip}
\newenvironment{definition}[2][]{\fkstmt[#1]{定义}{#2}}{\par\medskip}
\newenvironment{example}[2][]{\fkstmt[#1]{例}{#2}}{\par\medskip}
\newenvironment{remark}[2][]{\fkstmt[#1]{注}{#2}}{\par\medskip}
\newenvironment{notation}[2][]{\fkstmt[#1]{记号}{#2}}{\par\medskip}
\newenvironment{conjecture}[2][]{\fkstmt[#1]{猜想}{#2}}{\par\medskip}
\newenvironment{question}[2][]{\fkstmt[#1]{问题}{#2}}{\par\medskip}
\newenvironment{problem}[2][]{\fkstmt[#1]{问题}{#2}}{\par\medskip}
\newenvironment{assumption}[2][]{\fkstmt[#1]{假设}{#2}}{\par\medskip}
% 兼容房规旧习惯：带 p 后缀 = 证明紧跟（与原书无尾符的实情一致）
\newenvironment{theoremp}[2][]{\fkstmt[#1]{定理}{#2}}{\par\medskip}
\newenvironment{propositionp}[2][]{\fkstmt[#1]{命题}{#2}}{\par\medskip}
\newenvironment{lemmap}[2][]{\fkstmt[#1]{引理}{#2}}{\par\medskip}
\newenvironment{corollaryp}[2][]{\fkstmt[#1]{推论}{#2}}{\par\medskip}
\newenvironment{definitionp}[2][]{\fkstmt[#1]{定义}{#2}}{\par\medskip}
\newenvironment{examplep}[2][]{\fkstmt[#1]{例}{#2}}{\par\medskip}

% 证明：原书“Proof.”起首，末尾无符号（个别处有 “The proof is complete.” 之类的话）
\newenvironment{proof}[1][证明]{\par\medskip\noindent\textbf{#1}\ \ignorespaces}{\par\medskip}

% ---------- 译注（可见）----------
% 概念性批注用 \yizhu{…}：左侧细竖线 + 小字号，随文排，读者可见。
% 校勘性批注（笔误、页图与文字层不一致、术语取舍理由）用 LaTeX 注释
%   % [Tutor Note: …]  正文不印，由 work/tools/gen_notes.py 汇总进附录 B。
% 两类都进附录 B；见 STYLE_GUIDE.md §6。
\newtcolorbox{yizhubox}{
  enhanced, breakable, blanker,
  borderline west={1.4pt}{0pt}{gray!55},
  left=9pt, right=3pt, top=2pt, bottom=2pt,
}
\newcommand{\yizhu}[1]{\begin{yizhubox}\small\textbf{〔译注〕}\hspace{0.3em}#1\end{yizhubox}}

% ---------- 图 ----------
\newcommand{\figtitle}[2]{\par\smallskip\noindent{\centering\small 图~#1\quad #2\par}\smallskip}
\newcommand{\tabtitle}[2]{\par\smallskip\noindent{\centering\small 表~#1\quad #2\par}\smallskip}

% 交换图 / 示意图占位：\figdiag{原书页码}{该页第几个字母}
% 例：\figdiag{545}{a} → figures/dg_545_a.pdf（由 Lead 从 200dpi 页图裁切并核对）
% 文件未就位时印一个占位框，保证编译不中断。
\newcommand{\figdiag}[2]{%
  \begin{center}%
  \IfFileExists{figures/dg_#1_#2.pdf}%
    {\includegraphics[width=0.92\linewidth]{figures/dg_#1_#2.pdf}}%
    {\fbox{\parbox{0.55\linewidth}{\centering\small ［交换图：原文 p.#1（#2）待裁切］}}}%
  \end{center}}
% 单独排版的小图（原书无编号者）：
\newcommand{\figraw}[2]{%
  \begin{center}%
  \IfFileExists{figures/dg_#1_#2.pdf}%
    {\includegraphics[width=0.7\linewidth]{figures/dg_#1_#2.pdf}}%
    {\fbox{\parbox{0.45\linewidth}{\centering\small ［图：原文 p.#1（#2）待裁切］}}}%
  \end{center}}

% 内部 tikzcd 图（译者自己排的交换图；仅在原图简单到可无损重排时使用）
\newcommand{\diag}[1]{\begin{center}\begin{tikzcd}[row sep=2.2em,column sep=2.6em,ampersand replacement=\&]#1\end{tikzcd}\end{center}}

% ---------- 节首/章末（房规）----------
\newif\ifshowtakeaway     \showtakeawaytrue
\newif\ifshowchapterreview \showchapterreviewtrue
\newtcolorbox{takeawaybox}{
  enhanced, breakable, colback=blue!4!white, colframe=blue!55!black,
  boxrule=0.4pt, arc=1.5pt, left=5pt, right=5pt, top=3pt, bottom=3pt,
  parskip=0.4em,
  title={\small\bfseries 本节重点}, fonttitle=\small\bfseries,
  coltitle=white, colbacktitle=blue!55!black,
  attach boxed title to top left={xshift=4pt,yshift=-2pt},
}
\newtcolorbox{reviewbox}{
  enhanced, breakable, colback=orange!3!white, colframe=orange!65!black,
  boxrule=0.4pt, arc=1.5pt, left=6pt, right=6pt, top=4pt, bottom=4pt,
  parskip=0.6em,
  title={\small\bfseries 本章回顾与评注}, fonttitle=\small\bfseries,
  coltitle=white, colbacktitle=orange!65!black,
  attach boxed title to top left={xshift=4pt,yshift=-2pt},
}
\newcommand{\takeaway}[1]{\ifshowtakeaway\begin{takeawaybox}\small #1\end{takeawaybox}\fi}
\newcommand{\chapterreview}[1]{\ifshowchapterreview\begin{reviewbox}\small #1\end{reviewbox}\fi}

% ---------- 常用记号 ----------
\newcommand{\R}{\mathbb{R}}
\newcommand{\C}{\mathbb{C}}
\newcommand{\Z}{\mathbb{Z}}
\newcommand{\Q}{\mathbb{Q}}
\newcommand{\N}{\mathbb{N}}
\newcommand{\PP}{\mathbb{P}}
\newcommand{\bA}{\mathbb{A}}
\newcommand{\id}{\mathrm{id}}
\newcommand{\cO}{\mathcal{O}}
\newcommand{\cA}{\mathcal{A}}
\newcommand{\cB}{\mathcal{B}}
\newcommand{\cC}{\mathcal{C}}
\newcommand{\cD}{\mathcal{D}}
\newcommand{\cE}{\mathcal{E}}
\newcommand{\cF}{\mathcal{F}}
\newcommand{\cG}{\mathcal{G}}
\newcommand{\cH}{\mathcal{H}}
\newcommand{\cI}{\mathcal{I}}
\newcommand{\cJ}{\mathcal{J}}
\newcommand{\cK}{\mathcal{K}}
\newcommand{\cL}{\mathcal{L}}
\newcommand{\cM}{\mathcal{M}}
\newcommand{\cN}{\mathcal{N}}
\newcommand{\cP}{\mathcal{P}}
\newcommand{\cQ}{\mathcal{Q}}
\newcommand{\cR}{\mathcal{R}}
\newcommand{\cS}{\mathcal{S}}
\newcommand{\cT}{\mathcal{T}}
\newcommand{\cU}{\mathcal{U}}
\newcommand{\cV}{\mathcal{V}}
\newcommand{\cW}{\mathcal{W}}
\newcommand{\cX}{\mathcal{X}}
\newcommand{\cY}{\mathcal{Y}}
\newcommand{\cZ}{\mathcal{Z}}
\DeclareMathOperator{\Spec}{Spec}
\DeclareMathOperator{\Proj}{Proj}
\DeclareMathOperator{\Pic}{Pic}
\DeclareMathOperator{\Ext}{Ext}
\DeclareMathOperator{\Tor}{Tor}
\DeclareMathOperator{\RHom}{\mathbf{R}Hom}
\DeclareMathOperator{\Aut}{Aut}
\DeclareMathOperator{\Auteq}{Auteq}
\DeclareMathOperator{\Tot}{Tot}
\DeclareMathOperator{\Cone}{Cone}
\DeclareMathOperator{\Cyl}{Cyl}
\DeclareMathOperator{\ch}{ch}
\DeclareMathOperator{\td}{td}
\DeclareMathOperator{\Td}{Td}
\DeclareMathOperator{\rank}{rank}
\DeclareMathOperator{\im}{Im}
\DeclareMathOperator{\Ker}{Ker}
\DeclareMathOperator{\Coker}{Coker}
\DeclareMathOperator{\Supp}{Supp}
\DeclareMathOperator{\codim}{codim}
\DeclareMathOperator{\coker}{coker}
\DeclareMathOperator{\Ob}{Ob}
\DeclareMathOperator{\Mor}{Mor}
\DeclareMathOperator{\Hom}{Hom}
\DeclareMathOperator{\End}{End}
% 原书用粗体 D、R、L：\mathbf{D}^{b}(X)、\mathbf{R}f_*、\mathbf{L}f^{*}、\otimes^{\mathbf{L}}
\newcommand{\otL}{\otimes^{\mathbf{L}}}
\newcommand{\Rf}{\mathbf{R}}
\newcommand{\Lf}{\mathbf{L}}
\newcommand{\Rb}{\mathbf{R}^{\bullet}}
\newcommand{\Db}{\mathbf{D}^{b}}
\newcommand{\Dm}{\mathbf{D}^{-}}
\newcommand{\Dp}{\mathbf{D}^{+}}
\newcommand{\Du}{\mathbf{D}}
\newcommand{\db}{\mathbf{D}^{b}}
\newcommand{\dminus}{\mathbf{D}^{-}}
\newcommand{\dplus}{\mathbf{D}^{+}}
\newcommand{\dunb}{\mathbf{D}}
\newcommand{\coh}{\mathrm{coh}}
\newcommand{\Qcoh}{\mathrm{Qcoh}}
\newcommand{\Mod}{\mathrm{Mod}}
\newcommand{\Dcat}{\mathcal{D}}
\newcommand{\Tr}{\operatorname{Tr}}
\newcommand{\adj}{\dashv}
\newcommand{\iso}{\xrightarrow{\ \sim\ }}
\newcommand{\into}{\hookrightarrow}
\newcommand{\onto}{\twoheadrightarrow}
\newcommand{\sr}[1]{\mathrm{#1}}
\newcommand{\eps}{\varepsilon}
\newcommand{\vp}{\varphi}
\newcommand{\figfont}{\small}

% ---------- 参考文献（原书式：[编号] 或 [作者年]）----------
\newenvironment{biblist}{\begin{list}{}{\leftmargin=3.6em\itemindent=-3.6em
  \itemsep=0.6pt\parsep=0pt\topsep=2pt}}{\end{list}}
\newcommand{\bibentry}[2]{\item[\textnormal{[#1]}] #2\par}

% ---------- 术语表 / 索引 ----------
\newenvironment{termlist}{\begin{list}{}{\leftmargin=2.2em\itemindent=0pt
  \itemsep=0.4pt\parsep=0pt\topsep=2pt}}{\end{list}}
\newcommand{\termitem}[2]{\item #1\hfill{\small #2}}
\newenvironment{conceptlist}{\begin{list}{}{\leftmargin=0pt\itemindent=0pt
  \itemsep=2pt\parsep=0pt\topsep=2pt}}{\end{list}}
\newcommand{\concept}[4]{\item\textbf{#1}\quad{\small（#2）}\par\nopagebreak
  \hangindent=2em #3\par
  \hangindent=2em {\small\textbf{出处：}#4}\par}
% ---------- 分页 ----------
% openany 应让章首接排，但 ctexbook 仍会因 \cleardoublepage 产生空白页；
% 直接把它降级为 \clearpage，避免书里出现只带页眉的空白页。
\let\cleardoublepage\clearpage
